%\documentclass[tikz,border=10pt]{standalone}
%\usepackage{amsmath,amssymb}
%\usetikzlibrary{arrows.meta,positioning,calc,decorations.pathmorphing}
%
%\newcommand{\ord}{\operatorname{ord}}
%
%\begin{document}
%	\begin{tikzpicture}[
%		>=Latex,
%		font=\small,
%		panel/.style={
%			rounded corners=5pt, thick},
%		title/.style={font=\large\bfseries},
%		point/.style={circle, fill=black, inner sep=1.4pt},
%		good/.style={green!50!black, thick},
%		bad/.style={red!70!black, thick},
%		flow/.style={->, thick, black},
%		every node/.style={align=center}
%		]
%		
%		% =========================================================
%		% PANEL A: Local Correction (Sections 1 & 8)
%		% =========================================================
%		\node[panel, minimum width=6cm, minimum height=4.6cm] (A) at (0,0) {};
%		
%		% left local disk
%		\node[point,label={[yshift=-1mm]below:$p$}] at (-1.55,0.1) {};
%		\draw[dashed] (-1.55,0.1) circle (0.48);
%		\node[bad] at (-1.55,1.1) {$f=z^{-m}u$};
%		\node[red!75!black] at (-1.55,-1.08) {pole order $m$};
%		
%		% a little "blow-up" graphic
%		\draw[bad] (-1.92,-0.48) .. controls (-1.78,-0.05) and (-1.66,0.03) .. (-1.57,0.08);
%		\draw[bad] (-1.18,-0.48) .. controls (-1.32,-0.05) and (-1.44,0.03) .. (-1.53,0.08);
%		
%		% arrow
%		\draw[flow] (-0.52,0.1) -- (0.52,0.1);
%		\node[above] at (0,0.1) {$\times\, z^n$};
%		
%		% right local disk
%		\node[point,label={[yshift=-1mm]below:$p$}] at (1.55,0.1) {};
%		\draw[dashed] (1.55,0.1) circle (0.48);
%		\node[good] at (1.55,1.03) {$z^n f=z^{\,n-m}u$};
%		\draw[good] (1.15,-0.12) .. controls (1.34,0.30) and (1.76,0.32) .. (1.96,-0.08);
%		\node[blue!70!black,align=center] at (1.55,-1.08) {holomorphic\\ iff $n-m\ge 0$};
%	\end{tikzpicture}
%\end{document}


\documentclass[tikz,border=8pt]{standalone}
\usepackage{amsmath}
\usetikzlibrary{arrows.meta}

% Fonts
\usepackage[T1]{fontenc}
\usepackage[utf8]{inputenc}
\usepackage{newpxtext,newpxmath}
\usepackage{sectsty}

\renewcommand{\d}{\mathrm d}

\begin{document}
	
\begin{tikzpicture}[
	>=Latex,
	line cap=round,
	line join=round,
	font=\small,
	x=1cm,y=1cm
	]
	
	%------------------------------------------------
	% Parameters
	%------------------------------------------------
	\def\px{0.5}
	\def\py{-0.5}
	\def\rho{1}
	
	% Styles
	\tikzset{
		zero/.style={blue!70!black, fill=blue!70!black},
		polemark/.style={red!75!black, very thick},
		imgf/.style={violet, thick},
		imgg/.style={green!60!black, thick}
	}
	
	%====================================================
	% 1. Domain of f : z-plane
	%====================================================
	\begin{scope}[shift={(0,0)}]
		\node[font=\normalsize] at (0,-2.9) {$z$-plane (domain of $f$)};
		
		% axes
		\draw[->] (-2.8,0)--(2.8,0) node[right] {$\Re z$};
		\draw[->] (0,-2.4)--(0,2.8) node[above] {$\Im z$};
		
		% contour around p
		\draw[blue, thick] (\px,\py) circle (\rho);
		\draw[blue,->] (\px+\rho,\py) arc[start angle=0,end angle=35,radius=\rho];
		\node[blue] at (1.75,0.95) {$C_p:\ |z-p|=\rho$};
		
		% pole p
		\draw[polemark] (\px,\py) ++(-0.08,-0.08) -- ++(0.16,0.16);
		\draw[polemark] (\px,\py) ++(-0.08,0.08)  -- ++(0.16,-0.16);
		\node[red!75!black, right] at (\px+0.10,\py-0.16) {$p$};
%		\node[red!75!black] at (-1.45,-1.95) {pole of order $2$};
		
		% zero at z=-1
		\fill[zero] (-1,0) circle (2pt);
		\node[blue!70!black, above left] at (-1,0) {$-1$};
%		\node[blue!70!black] at (-1.55,1.65) {zero of order $1$};
	\end{scope}
	
	% Arrow for f
	\draw[->, red] (3.75,0.5) -- (7.75,0.5)
	node[midway, above] {$\displaystyle f(z)=\frac{z+1}{(z-p)^2}$}
	node[midway, below, align=center, red!70!black] {
		$\displaystyle \log f = \log (z+1)-2\log(z-p)$\\ [2pt]
		$\displaystyle \frac{\d f}{f}=\frac{\d z}{z+1}-2\,\frac{\d z}{z-p}$\\[4pt]
		$\displaystyle \operatorname{ord}_{-1}(f)=1,\; \operatorname{ord}_p(f)=-2$
	};
	
%	\draw[->, red!70!black] (3.2,-0.4) -- (7.2,-0.4)
%	node[midway, below, align=center] {%
%		$\displaystyle \operatorname{ord}_p(f)=-2,\qquad \operatorname{ord}_{-1}(f)=1$\\[2pt]
%		$\displaystyle \frac{df}{f}=\frac{dz}{z+1}-2\,\frac{dz}{z-p}$};
	
	%====================================================
	% 2. Codomain of f : f-plane
	%====================================================
	\begin{scope}[shift={(11,0)}]
		\node[font=\normalsize] at (0,-2.9) {$f(z)$-plane (codomain of $f$)};
		
		% axes
		\draw[->] (-2.8,0)--(2.8,0) node[right] {$\Re f$};
		\draw[->] (0,-2.4)--(0,2.8) node[above] {$\Im f$};
		
		% origin
		\fill (0,0) circle (1.6pt);
		\node[below right] at (0,0) {$0$};
		\node[blue!70!black, above left] at (0,0) {$f(-1)=0$};
		
		% image curve f(C_p)
		% z = p + rho e^{it}, f(z) = (z+1)/(z-p)^2
		\draw[imgf]
		plot[domain=0:6.283, samples=700]
		(
		{((\px+1+\rho*cos(\x r))*cos(2*\x r) + (\py+\rho*sin(\x r))*sin(2*\x r))/(\rho*\rho)},
		{((\py+\rho*sin(\x r))*cos(2*\x r) - (\px+1+\rho*cos(\x r))*sin(2*\x r))/(\rho*\rho)}
		);
		
		% manual arrows on the image curve
		\draw[imgf,->]
		({((\px+1+\rho*cos(0.8))*cos(1.6) + (\py+\rho*sin(0.8))*sin(1.6))/(\rho*\rho)},
		{((\py+\rho*sin(0.8))*cos(1.6) - (\px+1+\rho*cos(0.8))*sin(1.6))/(\rho*\rho)})
		--
		({((\px+1+\rho*cos(0.95))*cos(1.9) + (\py+\rho*sin(0.95))*sin(1.9))/(\rho*\rho)},
		{((\py+\rho*sin(0.95))*cos(1.9) - (\px+1+\rho*cos(0.95))*sin(1.9))/(\rho*\rho)});
		
		\draw[imgf,->]
		({((\px+1+\rho*cos(2.7))*cos(5.4) + (\py+\rho*sin(2.7))*sin(5.4))/(\rho*\rho)},
		{((\py+\rho*sin(2.7))*cos(5.4) - (\px+1+\rho*cos(2.7))*sin(5.4))/(\rho*\rho)})
		--
		({((\px+1+\rho*cos(2.85))*cos(5.7) + (\py+\rho*sin(2.85))*sin(5.7))/(\rho*\rho)},
		{((\py+\rho*sin(2.85))*cos(5.7) - (\px+1+\rho*cos(2.85))*sin(5.7))/(\rho*\rho)});
		
		\node[violet] at (2.1,2.0) {$f[C_p]$};
%		\node[align=center] at (0,-2.45) {twisted image curve\\due to the pole at $p$};
	\end{scope}
	
%%	\def\rho{1.75}
%	%====================================================
%	% 3. Domain of g : z-plane again
%	%====================================================
%	\begin{scope}[shift={(0,-6)}]
%		\node[font=\normalsize] at (0,-2.3) {$z$-plane (domain of $g$)};
%		
%		% axes
%		\draw[->] (-2.8,0)--(2.8,0) node[right] {$\Re z$};
%		\draw[->] (0,-1.8)--(0,1.8) node[above] {$\Im z$};
%		
%		% same contour around p
%		\draw[blue, thick] (\px,\py) circle (\rho);
%		\draw[blue,->] (\px+\rho,\py) arc[start angle=0,end angle=35,radius=\rho];
%		\node[blue] at (1.75,0.95) {$C_p:\ |z-p|=\rho$};
%		
%		% p regular for g
%		\fill[green!60!black] (\px,\py) circle (1.8pt);
%		\node[green!50!black, right] at (\px+0.10,\py-0.16) {$p$};
%%		\node[green!50!black] at (-1.5,-1.95) {regular point};
%		
%		% zero at z=-1
%		\fill[zero] (-1,0) circle (2pt);
%		\node[blue!70!black, above left] at (-1,0) {$-1$};
%%		\node[blue!70!black] at (-1.55,1.65) {zero of order $1$};
%	\end{scope}
%	
%	% Arrow for g
%	\draw[->, blue] (3.75,-5.5) -- (7.75,-5.5)
%	node[midway, above] {$\displaystyle g(z)=(z-p)^2f(z)=z+1$}
%	node[midway, below, align=center, blue!70!black] {%
%		$\displaystyle \log g = \log (z+1)$\\ [2pt]
%		$\displaystyle \frac{\d g}{g}=\frac{\d z}{z+1}$\\[4pt]
%		$\displaystyle \operatorname{ord}_{-1}(f)=1$
%%		$\displaystyle \operatorname{ord}_p(g)=0,\; \operatorname{ord}_{-1}(g)=1$\\[2pt]
%%		$\displaystyle \frac{df}{f}=\log $
%%%		$\displaystyle g \text{ is holomorphic at } p$
%	};
%	
%%	\draw[->, blue!70!black] (23.2,-0.4) -- (27.2,-0.4)
%%	node[midway, below, align=center] {%
%%		$\displaystyle \operatorname{ord}_p(g)=0,\qquad \operatorname{ord}_{-1}(g)=1$\\[2pt]
%%		$\displaystyle g \text{ is holomorphic at } p$};
%	
%	%====================================================
%	% 4. Codomain of g : g-plane
%	%====================================================
%	\begin{scope}[shift={(11,-6)}]
%		\node[font=\normalsize] at (0,-2.3) {$g(z)$-plane (codomain of $g$)};
%		
%		% axes
%		\draw[->] (-2.8,0)--(2.8,0) node[right] {$\Re g$};
%		\draw[->] (0,-1.8)--(0,1.8) node[above] {$\Im g$};
%		
%		% origin
%		\fill (0,0) circle (1.6pt);
%		\node[below left] at (0,0) {$0$};
%		\node[blue!70!black, above left] at (0,0) {$g(-1)=0$};
%		
%		% center g(p)=p+1
%		\fill[green!60!black] (\px+1,\py) circle (2pt);
%		\node[green!50!black, right] at (\px+1,\py) {$g(p)=p+1$};
%		
%		% image circle g(C_p)
%		\draw[imgg] ({\px+1},{\py}) circle (\rho);
%		\draw[imgg,->] ({\px+1+\rho},{\py}) arc[start angle=0,end angle=35,radius=\rho];
%		
%		\node[green!50!black] at (2.15,1) {$g[C_p]$};
%%		\node[align=center] at (0,-2.45) {ordinary translated circle\\no singular behavior};
%	\end{scope}
	
%	%====================================================
%	% Bottom summary
%	%====================================================
%	\node[draw=black, align=left, text width=18cm] at (8,-15) {%
%%		\texttt{Zeros and poles}\\[4pt]
%%		$f(z)=\dfrac{z+1}{(z-p)^2}$ has a zero at $z=-1$ and a pole of order $2$ at $z=p$.\\[6pt]
%%		$g(z)=(z-p)^2f(z)=z+1$ keeps the same zero at $z=-1$, but the pole at $p$ is canceled.\\[6pt]
%%		$\operatorname{ord}_{-1}(f)=\operatorname{ord}_{-1}(g)=1,\qquad
%%		\operatorname{ord}_p(f)=-2,\qquad
%%		\operatorname{ord}_p(g)=0.$%
%		$f$ has pole order at most 2 at $p$ $\iff$ $(z-p)^2f$ is holomorphic at $p$ $\iff$ $\operatorname{ord}_{p}((z-p)^2f)\geq 0$ $\iff$ $\operatorname{ord}_p(f)+2\geq 0$.
%	};
%
	\draw[dashed] (-3,-3.25) rectangle (15,3.5);
%	\draw[dashed] (-3,-8.75) rectangle (15,-3.5);
	
%	\draw[->,bend angle=45,start angle=0,end angle=90] (16, -1) to (16, -2);
	
\end{tikzpicture}
	
\end{document}